\documentclass[11pt,a4paper]{article}
% Common preamble for the hint sheets (\input from each file)
\usepackage[utf8]{inputenc}
\usepackage[T1]{fontenc}
\usepackage{lmodern}
\usepackage{amsmath,amssymb,amsthm,mathtools}
\usepackage[a4paper,margin=2.3cm]{geometry}
\usepackage{enumitem}
\usepackage{xcolor}
\usepackage[hidelinks]{hyperref}
\renewcommand{\leq}{\leqslant}
\renewcommand{\le}{\leqslant}
\renewcommand{\geq}{\geqslant}
\renewcommand{\ge}{\geqslant}
\setlength{\parindent}{0pt}
\setlength{\parskip}{0.35em}
% \hintstitle{Sheet title}{Class work / Homework}
\newcommand{\hintstitle}[2]{\begin{center}{\LARGE\bfseries Hints: #1\par}\vspace{0.5em}{\large #2\par}\end{center}
\noindent\textit{The hints to each problem are cumulative. Read one, try again, and open the next only if you are still stuck.}\par\medskip\hrule\medskip}
\newcommand{\problem}[1]{\par\bigskip{\large\bfseries Problem #1}\par\nopagebreak\smallskip}
\newcommand{\subpart}[1]{\par\smallskip\textbf{(#1)}\par\nopagebreak}
\newcommand{\hint}[1]{\par\smallskip\noindent\textcolor{black!70}{\textbf{Hint #1.}}\enspace}
\newcommand{\note}[1]{\par\smallskip\noindent\textit{Note.} #1\par}
\providecommand{\lt}{<}
\providecommand{\gt}{>}
\providecommand{\tightlist}{\setlength{\itemsep}{0pt}\setlength{\parskip}{0pt}}
\emergencystretch=3em

\begin{document}
\hintstitle{Integrals -- 5. The gamma and beta functions and Wallis' formulas}{Class work}

\problem{1}
\hint{1}
Write \(\Gamma(p)\Gamma(q)=\int_0^\infty\int_0^\infty x^{p-1}y^{q-1}e^{-x-y}\,dx\,dy\) and substitute \(x=u^2\), \(y=v^2\): it becomes \(4\iint_{u,v\gt0}u^{2p-1}v^{2q-1}e^{-(u^2+v^2)}du\,dv\).
\hint{2}
In polar coordinates \(u=r\cos\theta\), \(v=r\sin\theta\) the integral factors: \(2\int_0^\infty r^{2(p+q)-1}e^{-r^2}dr=\Gamma(p+q)\) and \(2\int_0^{\pi/2}\cos^{2p-1}\theta\sin^{2q-1}\theta\,d\theta=B(p,q)\) (Fact 1). For \(p=q=\frac12\): \(\Gamma(\frac12)^2=\Gamma(1)B(\frac12,\frac12)=2\cdot\frac\pi2=\pi\).

\problem{2}
\hint{1}
Theorem 2 is Fact 1 with \(2p-1=n\), \(2q-1=0\). For even \(n=2m\) use \(\Gamma(m+\frac12)=\frac{(2m-1)!!}{2^m}\sqrt\pi\) to get \(\int_0^{\pi/2}\sin^{2m}x\,dx=\frac{(2m-1)!!}{(2m)!!}\cdot\frac\pi2\); for odd \(n=2m+1\) you get \(\frac{(2m)!!}{(2m+1)!!}\).
\hint{2}
For the product: \(I_n=\int_0^{\pi/2}\sin^nx\,dx\) decreases in \(n\), so \(I_{2m+1}\le I_{2m}\le I_{2m-1}\), and \(\frac{I_{2m-1}}{I_{2m+1}}=\frac{2m+1}{2m}\to1\); hence \(\frac{I_{2m}}{I_{2m+1}}\to1\), which is the Wallis product.
\hint{3}
\(\int_0^{2\pi}\cos^{2020}x\,dx=4I_{2020}=2\pi\binom{2020}{1010}2^{-2020}\). For the second, \(\sin^2x\cos^2x=\frac14\sin^22x\), so the integral is \(\frac\pi8\).

\problem{3}
\hint{1}
Bring each integral to the form \(\int_0^1u^{p-1}(1-u)^{q-1}du\): (a) is already \(B(\frac32,\frac32)\); (b) put \(u=x^2\); (c) put \(x=t^2\); (d) put \(u=x^{2025}\).
\hint{2}
\begin{enumerate}
\def\labelenumi{(\alph{enumi})}
\tightlist
\item
  \(\frac{\Gamma(3/2)^2}{\Gamma(3)}=\frac\pi8\); (b) \(\frac12B(\frac12,\frac32)=\frac\pi4\); (c) \(2B(2,\frac32)=\frac8{15}\); (d) \(\frac1{2025}B(1,2026)=\frac1{2025\cdot2026}\).
\end{enumerate}

\problem{4}
\hint{1}
Substitute \(x=\frac{t}{1-t}\) (the inverse of \(t=\frac x{1+x}\)): then \(1+x=\frac1{1-t}\), \(dx=\frac{dt}{(1-t)^2}\) and the integral becomes \(\int_0^1t^{a-1}(1-t)^{-a}dt=B(a,1-a)\).
\hint{2}
\(B(a,1-a)=\Gamma(a)\Gamma(1-a)=\frac\pi{\sin\pi a}\). The contour method gets the same value from the residue at \(-1\) of \(\frac{z^{a-1}}{1+z}\) on a keyhole contour; here the work is hidden in the reflection formula (which in turn can be proved by this very integral).

\problem{5}
\hint{1}
Pointwise, \(e^{-x^{2n}}\to1\) for \(\lvert x\rvert\lt1\) and \(\to0\) for \(\lvert x\rvert\gt1\); dominate by \(e^{-x^2}\) for \(\lvert x\rvert\ge1\) and by \(1\) for \(\lvert x\rvert\lt1\). The graph tends to the indicator of \([-1,1]\), a rectangle of area \(2\).
\hint{2}
Exactly: \(u=x^{2n}\) gives \(\int_{-\infty}^\infty e^{-x^{2n}}dx=\frac1n\Gamma\left(\frac1{2n}\right)=2\Gamma\left(1+\frac1{2n}\right)\to2\Gamma(1)=2\).
\end{document}
